Posons $s = a + b + c + d$.

\textbf{1.} Chaque dénominateur est obtenu en retirant de $s$ un nombre
strictement positif, donc est strictement inférieur à $s$. Ainsi
\[
S > \frac{a}{s} + \frac{b}{s} + \frac{c}{s} + \frac{d}{s} = 1.
\]

\textbf{2.} On a $a + b + d > a + b$ et $a + b + c > a + b$, donc
\[
\frac{a}{a+b+d} + \frac{b}{a+b+c} < \frac{a}{a+b} + \frac{b}{a+b} = 1.
\]
De même, $b + c + d > c + d$ et $a + c + d > c + d$ donnent
$\frac{c}{b+c+d} + \frac{d}{a+c+d} < \frac{c}{c+d} + \frac{d}{c+d} = 1$.
En additionnant, $S < 2$.

\textbf{3.} Pour $t > 0$ :
\[
\begin{aligned}
S(t, t, 1, 1) &= \frac{t}{2t+1} + \frac{t}{2t+1} + \frac{1}{t+2} + \frac{1}{t+2}
= \frac{2t}{2t+1} + \frac{2}{t+2}, \\
S(1, t, 1, t) &= \frac{1}{1+2t} + \frac{t}{2+t} + \frac{1}{1+2t} + \frac{t}{2+t}
= \frac{2}{1+2t} + \frac{2t}{2+t}.
\end{aligned}
\]
Quand $t \to 0^+$, $S(t, t, 1, 1) \to 0 + 1 = 1$ et $S(1, t, 1, t) \to 2 + 0 = 2$.

\textbf{4.} Soit $\sigma \in \left]1, 2\right[$. D'après la question 3, il
existe $t_1 > 0$ tel que $S(t_1, t_1, 1, 1) < \sigma$ et $t_2 > 0$ tel que
$S(1, t_2, 1, t_2) > \sigma$. Posons $P = (t_1, t_1, 1, 1)$ et
$Q = (1, t_2, 1, t_2)$. Pour $\lambda \in [0, 1]$, chaque coordonnée de
$(1-\lambda)P + \lambda Q$ est une moyenne pondérée de deux nombres
strictement positifs, donc est strictement positive : $g$ est bien définie
sur $[0, 1]$. Elle est continue, comme somme de quotients de fonctions
affines de $\lambda$ dont les dénominateurs ne s'annulent pas. Comme
$g(0) = S(P) < \sigma < S(Q) = g(1)$, le théorème des valeurs
intermédiaires fournit $\lambda_0 \in \left]0, 1\right[$ tel que
$g(\lambda_0) = \sigma$.

\textbf{5.} D'après les questions 1 et 2, toutes les valeurs de $S$ sont
dans $\left]1, 2\right[$ ; d'après la question 4, chaque réel de
$\left]1, 2\right[$ est atteint. L'ensemble des valeurs prises par $S$ est
donc exactement l'intervalle ouvert $\left]1, 2\right[$. \hfill $\blacksquare$
